\subsection{Every rank-two source representation has conductor three}
\label{sec:gl2-conductor}

\begin{lemma}[K17-GL2-CONDUCTOR]\label{lem:gl2-conductor}
Let $e\geq1$, $\ell>3$, and $E=\mathbf Q(\sqrt{-3})$.
Every representation in the actual source class
$\mathrm{Adm2A}(\iota,e,6)$ has, at a split prime $p\ne3,\ell$,
the trace
\[
 \iota\operatorname{tr}\rho(\sigma^{-1})
       =1+\frac{a_p(f)}{p^e},
 \qquad
 f\in S_{2e+1}(\Gamma_0(3),\chi_{-3}),
\]
where $f$ is a normalized cuspidal eigenform and $\sigma$ is arithmetic
Frobenius at the selected place of $E$ above $p$.
In fact the primitive conductor of $f$ is exactly three.
\end{lemma}

\begin{proof}
The definition in
\texttt{AutomorphyDescent/AutNewformW.lean} supplies a continuous
irreducible $W:G_{\mathbf Q}\to\mathrm{GL}_2(\overline{\mathbf Q}_\ell)$,
a weight-$2e+1$ newform attached to $R=W\chi_\ell^{-e}$,
\[
 \mathrm{Level}(\ell,6,(e,-e),W),
 \qquad \rho\simeq\mathbf1\oplus W|_{G_E}.
 \tag{GC.1}
\]
We retain this literal level condition. In characteristic zero, fixing
all sixth roots of unity is equivalent to fixing all third roots of
unity. Thus every character satisfying the source's
\texttt{KillsMu 6} factors through
$\operatorname{Gal}(\mathbf Q(\mu_3)/\mathbf Q)$ and is either
$1$ or $\chi_{-3}$.

The source attachment in \texttt{AutomorphicForms/Newform.lean}
asserts semisimplicity and equality of the characteristic polynomials
at all but finitely many geometric Frobenius elements. Chebotarev
density and Brauer--Nesbitt identify $R$ with the Deligne representation
of the corresponding primitive eigenform $f$. The primitive level divides
the level in the source definition, and its explicit minimal-level clause
excludes a proper divisor. We may therefore use the usual primitive
newform and its finite coefficient field. The determinant of $R$ is
$\epsilon\chi_\ell^{-2e}$, where $\epsilon$ is the finite nebentypus.
The odd weight gives $\epsilon(-1)=-1$.

At a prime $q\ne\ell$, (GC.1) diagonalizes the actual inertia action of
$W$ with characters in $\{1,\chi_{-3}\}$. The cyclotomic twist is
unramified there, so the same holds for $R$. The finite inertia action
has zero monodromy: its logarithm on an open inertia subgroup is zero.
Consequently its Weil--Deligne representation is unramified at
$q\ne3,\ell$; at three it has zero monodromy and two tame characters,
each $1$ or $\chi_{-3}$. Carayol
\cite[\S0.7, Theorem~A, pp.~410--411]{K17Carayol1986}
identifies these local parameters with those of $f$. The additional
discrete-series hypothesis in his \S0.3 is required only when the
totally real field has even degree, so it is absent over $\mathbf Q$.

We give the coefficient-prime argument explicitly. Saito's
local compatibility theorem, in the exact author-stated form
\cite[\S2c, Theorem~3, p.~9; definitions on pp.~7--8]{K17SaitoSurvey},
gives potential semistability and identifies the Frobenius-semisimplified
\emph{full} Weil--Deligne representation, including monodromy, with the
local Langlands parameter at $\ell$. The theorem concerns classical
modular forms over $\mathbf Q$ and cites his 1997 paper; we are not
specializing a Hilbert theorem whose standing hypothesis excludes degree
one.

The flag in (GC.1), restricted to $G_{\mathbf Q_\ell}$ and twisted,
has characters
\[
       \delta_1,\qquad \delta_2\chi_\ell^{-2e},
 \tag{GC.2}
\]
where the $\delta_i$ are unramified and continuous, since
$\chi_{-3}$ is unramified at $\ell>3$.
Their Frobenius values are $\ell$-adic units: the image of the compact
group $G_{\mathbf Q_\ell}$ is compact, and its valuation image is a
finite subgroup of the additive group of rational numbers, hence zero.
An unramified character is crystalline with slope zero; a cyclotomic
power shifts that slope by its integer exponent, up to the common
choice of Frobenius convention. Thus the two characters in (GC.2)
have trivial Weil--Deligne inertia and slopes differing by $2e$.
These rank-one assertions also follow directly from the period $t$
and the one-dimensional calculation in
\cite[\S4.2, Proposition~4.4 and its proof, p.~24]{K17ColmezFontaine2000}.

Choose a finite local extension over which $R$ is semistable and a
finite coefficient extension containing its flag. The exactness of
$D_{\mathrm{st}}$ on semistable representations
\cite[\S4.1, Proposition~4.2, p.~23]{K17ColmezFontaine2000}
shows that the two Frobenius slopes of $D_{\mathrm{pst}}(R)$ are exactly
the two slopes in (GC.2). On a residue-degree-$f_0$ extension, use
$\varphi^{f_0}$ and divide valuations by $f_0$; the same difference is
$2e$. The relation
\[
                  N\varphi=\ell\varphi N
\]
is the defining relation of a $(\varphi,N)$-module
\cite[\S3.3, p.~21]{K17ColmezFontaine2000}.
If $N\ne0$ in dimension two, it maps one Frobenius generalized
eigenspace to the other and forces their normalized slopes to differ
by one. Since $2e\ge2$, this is impossible. Thus $N=0$.
Weil--Deligne inertia has finite image and its two graded characters
are trivial. Semisimplicity of finite-group representations in
characteristic zero therefore makes the full inertia action trivial.
Saito's theorem now says that $f$ is unramified at $\ell$.
Frobenius semisimplification preserves both inertia and the zero
monodromy assertion used here.

It follows that $\epsilon$ is unramified outside three, and its inertia
at three is tame of order at most two. Its primitive Dirichlet conductor
therefore divides three. Since it is odd, $\epsilon=\chi_{-3}$.
The determinant of the two inertia characters at three is accordingly
$\chi_{-3}$. Exactly one of them is nontrivial. Their tame Artin
conductor is one, and zero monodromy contributes nothing further.
Local compatibility and the conductor-preserving local Langlands
correspondence give primitive level exactly three. This proves the
assertion about the complete classical modular-form space.

Finally, $\chi_\ell(\sigma)=p$ and the source uses geometric Frobenius,
so
\[
 R(\sigma^{-1})=p^eW(\sigma^{-1}),\qquad
 \iota\operatorname{tr}R(\sigma^{-1})=a_p(f).
\]
Together with (GC.1), these equalities give the stated trace formula.
\end{proof}

\paragraph{Geometric meaning of the local step.}
The representation of $f$ is cut out of the etale cohomology of a
Kuga--Sato variety. At the coefficient prime its potentially semistable
realization is a log-crystalline $(\varphi,N)$-module. The separated
ordinary weights force $N=0$; the finite descent inertia is then trivial.
This is the precise crystalline input behind the conductor reduction.
The use of Saito's theorem and of $D_{\mathrm{st}}$ exactness is a
recheck of cited results, not a new proof of the comparison theorem.
